\subsection{Motivation}

Within the computer science program at DHBW, the lectures ``Digitaltechnik'' and ``Rechnerarchitektur'' teach important fundamentals for an understanding of modern digital computing systems. Both lectures can benefit from a good visualisation of the taught concepts. ``Digitaltechnik'' imparts fundamentals of digital circuits, while ``Rechnerarchitektur'' shows the internals of digital processors on a higher level of abstraction. An illustrative model of a simple \ac{cpu} could serve as a helpful bridge between the two levels of abstraction. From the perspective of ``Digitaltechnik'', individual modules of the \ac{cpu} are especially interesting. For example, registers or individual arithmetic operations can be shown on their physical counterparts. From the perspective of ``Rechnerarchitektur'', the composition and interconnection of the modules is more relevant. Here, the datapath or the control logic can be visualized. Both lectures may benefit from a visualisation of the order of steps within a digital system. Since the implementation is on a gate level, minimizing the number of gates is the most important metric when designing an implementation of the \ac{isa}.

This thought forms the core of the motivation for this work. Additionally, the assignment offers an enormous potential for learning the student. It touches a broad spectrum of technical topics, whose fundamentals were already taught. Therefore, these foundations can be strengthened and expanded upon. Planing and executing a complex technical project independently offers further opportunities for personal development.